\begin{proposition}\label{prop:r3472-maximal-third-indicator}
In the canonical positive spherical structure of every pseudo-unitary
categorification of $\sha{R8_3_3472}$,
the third Frobenius--Schur indicator of every nonunit simple is $2$.
Equivalently, rotation acts identically on each two-dimensional space
$\operatorname{Hom}(1,X_i^{\otimes3})$, $i>0$.
\end{proposition}
\begin{proof}
Put $\delta=7+5\sqrt2$, $D=700+490\sqrt2$, and $D'=700-490\sqrt2$.
The source dimensions are
$(1,\delta-1,\delta-1,\delta,\delta-1,\delta+1,\delta+1,\delta+1)$,
and $\dim\operatorname{Hom}(1,X_i^{\otimes3})=2$ for all $i>0$.
The spherical Cauchy theorem~\cite[Theorem~3.9]{BNRWCauchy} and
$DD'=9800$ imply that $3$ is coprime to
$N=\operatorname{FSexp}(\mathcal C)$.
The dimensions lie in $\mathbb Q(\zeta_N)$, so $8\mid N$.
For $3\widetilde3\equiv1\pmod N$, the automorphism
$\sigma(\zeta_N)=\zeta_N^{\widetilde3}$ therefore flips $\sqrt2$.
The higher Gauss-sum formula~\cite[Theorem~4.4, (4.17)]{NSWHigherGauss}
gives
\[
 \sum_i d_i\nu_3(X_i)=\frac D{\sigma(D)}
 =\frac D{D'}=\delta^2=1+2\sum_{i>0}d_i.
\]
We have $\nu_3(1)=1$. For $i>0$, $\nu_3(X_i)$ is the trace of an
order-three operator on a two-dimensional space, so
$\operatorname{Re}\nu_3(X_i)\le2$.
All $d_i$ are positive, so the displayed equality forces equality
for every summand. Every rotation eigenvalue is consequently $1$;
the operator is diagonalizable because its cube is the identity.
\end{proof}

The existing height-four full-character-field phase witness already
has third-rotation weights $(0,2)$, for trace types $(-1/2,1)$,
at every nonunit simple. Thus this restriction does not exclude that
relaxation. Exact source arithmetic and the cached-coordinate audit are
recorded in
\path{evidence/r3472_focus_20261022/independent/verify_third_indicators.py}.

% Bibliography addition needed if this fragment is integrated:
% \bibitem{NSWHigherGauss} S.-H.~Ng, A.~Schopieray and Y.~Wang,
% \emph{Higher Gauss sums of modular categories},
% Selecta Math. (N.S.) \textbf{25} (2019), Paper 53;
% arXiv:1812.11234v2, Theorem~4.4.
